\documentclass[answers,12pt,addpoints]{exam} 
\usepackage{import}

\import{C:/Users/prana/OneDrive/Desktop/MathNotes}{style.tex}

% Header
\newcommand{\name}{Pranav Tikkawar}
\newcommand{\course}{Spatial Statistics}
\newcommand{\assignment}{Homework 3.1}
\author{\name}
\title{\course \ - \assignment}

\begin{document}
\maketitle
\begin{questions}
\question I was curious if there was a method that could interpolate between two correlation functions. If it were possible, then we could use it to fulfil the property of the isotropy in local areas but some global anisotropy. I cant real think of any method to do this in a general case, but if this would be a reasonable thing to do.

\question Is there a measure of how much a correlation function is anisotropic? I would assume something relating to the eigenvalues of the covariance matrix. Where the ratio of the largest to smallest eigenvalue would be a measure of how anisotropic the correlation function is. As well as the eigenvectors would give us the direction of the anisotropy. 

\question I didn't consider that different parametrization of the same model could be difficult to interpret. It is clear for isotopic models what the form of the model should take, but for anisotropic models, it is not clear how to interpret the parameters just by looking at the model. 

\question I think the idea of using the same model but different forms with different parameters is interesting. The scale that the model is viewed on definitely can make a difference in interpretation of the model. Would it be reasonable to have different parameters for different scales even though they are the same model? I think it would be interesting to see how the parameters change as the scale changes.

\question I have recently starting learning more about the field of reproducing kernel Hilbert spaces and I feel like there is a natural application of that theory here, but I unfortunately I don't know enough about the field at the moment to make any concrete connections. I am going based on the intuition that for these kind of optimization/regression problems there is a natural way to create a canonical coordinate system with corresponding kernel that would make the problem easier to solve. I am not sure if this is the case, but it would be interesting to see if there is a connection between the two fields.

\question This weeks reading was readily short on theory which is nice because I feel like I got a good amount of examples for these anisotropic models. 
\end{questions}

\end{document}